\documentclass[letterpaper,11pt]{article}
\usepackage[top=0.8in, bottom=0.8in, left=0.9in, right=0.9in]{geometry}
\usepackage[hidelinks]{hyperref}
\usepackage{lmodern}
\usepackage{parskip}
\usepackage{microtype}
\setlength{\parskip}{6pt}
\setlength{\parindent}{0pt}
\begin{document}
\begin{flushleft}
\textbf{\large Ajay Venkatesh} \\
Brooklyn, NY \\
+1 516 585 9013 \\
\href{mailto:av3855@nyu.edu}{av3855@nyu.edu}
\end{flushleft}
\vspace{10pt}
June 8, 2026
\vspace{6pt}

Hiring Manager \\
Optimal Dynamics

\vspace{10pt}

Dear Hiring Manager,

I'm Ajay Venkatesh, completing an M.S. in Computer Engineering at NYU, with production software experience at PwC in Bengaluru, a summer SDE internship at Amazon, and ongoing on-campus work at NYU's Novel AI Technologies. I'm applying for the OR Software Engineer role on your OR/AI team, where I'd contribute \textbf{Python}-based optimization and machine learning software that supports logistics decision-making at scale.

At Amazon, I built \textbf{LLM-based retrieval pipelines (RAG)} with agentic workflows and an \textbf{MCP server}, and engineered backend services with \textbf{OOP design patterns}, test coverage, and \textbf{CI/CD} releases on \textbf{AWS}. That work maps closely to researching, implementing, and maintaining ML infrastructure in production with an emphasis on stability and measurable quality.

In my \textbf{RoBERTa LoRA fine-tuning} project, I ran hyperparameter sweeps and rank/learning-rate experiments under strict memory constraints, applying \textbf{PyTorch}, \textbf{Hugging Face}, and \textbf{PEFT} to improve model accuracy with minimal trainable parameters. At NYU and CampusMesh, I shipped \textbf{Python ETL pipelines}, cloud-hosted \textbf{AI/ML models}, and \textbf{embedding-based semantic search} systems backed by \textbf{PostgreSQL}, \textbf{AWS}, and \textbf{Docker}.

My stack includes \textbf{Python}, \textbf{Java}, \textbf{C++}, \textbf{PyTorch}, \textbf{AWS}, \textbf{Azure}, \textbf{Docker}, \textbf{Git}, \textbf{Pandas}, \textbf{NumPy}, \textbf{Cursor}, and \textbf{GitHub Copilot}. \textbf{Stochastic optimization} and dedicated \textbf{supply chain logistics} experience are areas I have not owned in production yet, but my coursework in \textbf{machine learning}, \textbf{deep learning}, and \textbf{big data} plus hands-on model tuning give me a strong foundation to ramp on your OR tooling.

What pulls me toward Optimal Dynamics is the chance to build decision software where optimization and AI change how freight and logistics operate in the real world. I care about rigorous engineering paired with learning systems that compound in production, and I'd welcome contributing to that mission on your OR/AI team.

I'd welcome the chance to discuss further. Thanks for considering my application.

\vspace{12pt}
Sincerely, \\
Ajay Venkatesh
\end{document}
